\begin{theorem}[Scalar calculus of character projectors]
    \label{thm:peps_weighted_character_projectors}
    \lean{Representation.sum_smul_charProjector_apply_of_mem,
        Representation.sum_smul_charProjector_commute,
        Representation.charProjector_mul_self}
    \leanok
    \uses{thm:peps_isotypic_projectors}
    Let $\rho$ be a finite-dimensional complex representation of a finite
    group. Given complex weights $f(\chi)$ on the occurring irreducible
    characters, put $A_f=\sum_\chi f(\chi)P_\chi$.
    If $v$ lies in an irreducible summand with character $\psi$, then
    $A_fv=f(\psi)v$. Moreover, $A_f\rho(g)=\rho(g)A_f$.
    Each projector $P_\chi$ is idempotent.
    These formulas express the isotypic decomposition used in
    \cite[Lemma~4.4]{Schuch2010PEPS}.
\end{theorem}
\begin{proof}\leanok
    On the summand with character $\psi$, the projector $P_\chi$ is the
    identity if $\chi=\psi$ and zero otherwise. Thus only its recorded
    weight survives. Each summand is preserved by the group action,
    which commutes with scalar multiplication. The irreducible summands
    span the whole representation space.
\end{proof}

\begin{definition}[Reciprocal weighted trace operator]
    \label{def:peps_delta_inverse}
    \lean{Representation.deltaOperatorInverse}
    \leanok
    \uses{def:peps_delta_operator}
    For the occurring irreducible characters, with dimensions $d_\chi$
    and multiplicities $m_\chi$, define
    \begin{align}
        \Delta^{-1}=\sum_\chi\frac{|G|m_\chi}{d_\chi}P_\chi.
        \label{eq:peps_delta_inverse}
    \end{align}
    These are the reciprocal coefficients of the weighted trace operator
    in \cite[Lemma~4.4]{Schuch2010PEPS}.
\end{definition}

\begin{theorem}[Invertibility of the weighted trace operator]
    \label{thm:peps_delta_inverse}
    \lean{Representation.deltaOperatorInverse_apply_of_mem,
        Representation.deltaOperatorInverse_mul,
        Representation.deltaOperator_mul_inverse,
        Representation.isUnit_deltaOperator,
        Representation.deltaOperatorInverse_commute,
        Representation.deltaOperator_commute}
    \leanok
    \uses{def:peps_delta_inverse, thm:peps_weighted_character_projectors,
        thm:peps_delta_expansion}
    The operator in (\ref{eq:peps_delta_inverse}) is the two-sided inverse
    of $\Delta$. On a summand with character $\chi$, it acts by
    $|G|m_\chi/d_\chi$. Both $\Delta$ and its inverse commute with every
    group operator. These assertions hold without semi-regularity or
    unitarity, for every finite-dimensional complex representation of a
    finite group.
\end{theorem}
\begin{proof}\leanok
    The dimensions and multiplicities of the occurring summands are
    positive. On each irreducible summand the two coefficients multiply
    to one:
    $[d_\chi/(|G|m_\chi)][|G|m_\chi/d_\chi]=1$.
    Hence $\Delta^{-1}\Delta=I$ on their direct sum. In finite dimension
    a left inverse is also a right inverse. Commutation follows from
    the scalar calculus of the character projectors.
\end{proof}

\begin{theorem}[Positive weighted trace operator]
    \label{thm:peps_delta_positive}
    \lean{Representation.adjoint_subrepresentation_of_unitary,
        Representation.charProjector_adjoint_of_unitary,
        Representation.isPositive_charProjector_of_unitary,
        Representation.isPositive_sum_smul_charProjector_of_nonneg,
        Representation.isPositive_deltaOperator_of_unitary,
        Representation.posDef_toMatrix_deltaOperator_of_unitary}
    \leanok
    \uses{thm:peps_delta_inverse, thm:peps_isotypic_projectors}
    Let $\rho$ be a unitary finite-dimensional representation of a finite
    group. Every invariant subrepresentation is unitary. Then $\Delta$ is positive, and its matrix in every orthonormal
    basis is positive definite. For unitary representations, the character
    projector of any irreducible summand is self-adjoint and positive.
    More generally, every nonnegative weighted sum of the occurring
    character projectors is positive. This supplies the positive bond weights
    appearing in \cite[Lemma~4.4]{Schuch2010PEPS}.
\end{theorem}
\begin{proof}\leanok
    Unitarity gives $\chi(g^{-1})=\overline{\chi(g)}$. Hence the character
    projectors of the occurring summands are self-adjoint. They are also
    idempotent, so they are positive. The coefficients
    $d_\chi/(|G|m_\chi)$ are positive. Their weighted sum is therefore
    positive; its invertibility makes it positive definite.
\end{proof}
